% ============================================================
% UGA-Themed Beamer Template (Clean MWE)
% ============================================================

\documentclass[aspectratio=169,professionalfonts]{beamer}

% ---------- COLORS (swap with official UGA hex if you like) ----------
\usepackage{xcolor}
\definecolor{UGARed}{HTML}{BA0C2F}   % Verify with brand guide
\definecolor{UGABlack}{HTML}{000000}
\definecolor{UGAGray}{HTML}{8A8D8F}

% ---------- THEME & PROGRESS DOTS ----------
\usetheme{default}
\useinnertheme{circles}
\useoutertheme[subsection=false]{miniframes}
\usefonttheme{professionalfonts}
\setbeamertemplate{navigation symbols}{}
% Head/foot colors
\setbeamercolor{subsection in head/foot}{fg=white,bg=UGARed}
\setbeamercolor{headline}{fg=white, bg=UGARed}

\setbeamercolor{section in head/foot}{fg=white,bg=UGARed}
\setbeamercolor{section in head/foot shaded}{fg=white,bg=UGARed}
\setbeamercolor{footline}{bg=UGABlack,fg=white}
\setbeamercolor{frametitle}{fg=UGABlack}
\setbeamercolor{title}{fg=UGABlack}
\setbeamercolor{normal text}{fg=UGABlack}
\setbeamercolor{structure}{fg=UGARed}

% --- HEADLINE: slide dots grouped by section ---
\setbeamersize{mini frame size=0.16cm,mini frame offset=0pt}
% Show section labels above the clickable slide dots.
\setbeamertemplate{section in head/foot}[default]
\setbeamertemplate{section in head/foot shaded}[default][75]
\setbeamertemplate{mini frame}{%
  \tikz[baseline=-0.5ex] \filldraw[draw=white,fill=white,line width=0.6pt] (0,0) circle (2pt);%
}
\setbeamertemplate{mini frame in current subsection}{%
  \tikz[baseline=-0.5ex] \filldraw[draw=white,fill=UGARed,line width=0.6pt] (0,0) circle (2pt);%
}
\setbeamertemplate{mini frame in other subsection}{%
  \usebeamertemplate{mini frame in current subsection}%
}
\setbeamertemplate{mini frame in other section}{%
  \usebeamertemplate{mini frame in current subsection}%
}
\setbeamertemplate{headline}{%
  \begin{beamercolorbox}[wd=\paperwidth]{headline}
    \vskip1ex
    \insertnavigation{\paperwidth}
    \vskip3ex
  \end{beamercolorbox}
}

% Simple footline with slide numbers
\setbeamertemplate{footline}{
  \leavevmode%
  \hbox{%
    \begin{beamercolorbox}[wd=\paperwidth,ht=2.5ex,dp=1ex,leftskip=1em,rightskip=1em]{footline}%
      \insertshortauthor{} \hfill \insertframenumber/\inserttotalframenumber
    \end{beamercolorbox}}%
}

% ---------- PACKAGES ----------
\usepackage{amsmath,amssymb,mathtools}
\usepackage{graphicx}
\usepackage{booktabs}
\usepackage{bbm}
\usepackage{array}
\usepackage{adjustbox}
\usepackage{hyperref}
\usepackage{tikz}
\usetikzlibrary{arrows.meta, positioning}
\usepackage{tcolorbox}
\usepackage{eurosym}
\usepackage{threeparttable}
% ---------- METADATA (fill these later) ----------
\title{Optimal Dynamic Pricing With Preferences for Novelty}
\subtitle{}
\author{Tate Mason}
\institute{The University of Georgia}
\date{\today}

% ---------- MACROS ----------
% Numbered equation with label
\newcommand{\beq}[1]{\begin{equation}\label{#1}}
\newcommand{\eeq}{\end{equation}}
% Accent color highlight
\newcommand{\hl}[1]{\textbf{\textcolor{UGARed}{#1}}}
% Simple figure placeholder macro
\newcommand{\figph}[2][0.9\linewidth]{%
  \begin{figure}
    \centering
    \fbox{\rule{0pt}{0.55#1}\rule{#1}{0pt}}%
    \caption{#2}
  \end{figure}
}

% ============================================================

% Keep navigation links in the theme's white foreground on the red header.
\hypersetup{colorlinks=true, linkcolor=., urlcolor=UGARed}


% Short-version navigation only. Original frame content is preserved verbatim.
\newcommand{\shortbuttons}{}
\newcommand{\shortmaincount}{27}
\IfFileExists{assets/quantified_theory/plain/optimal_prices_by_gamma.pdf}{\renewcommand{\shortmaincount}{33}}{}
\newif\ifshortappendix
\shortappendixfalse
\setbeamertemplate{footline}{%
  \leavevmode\vbox{\hbox{%
  \begin{beamercolorbox}[wd=\paperwidth,ht=3.8ex,dp=.6ex,rightskip=1em]{normal text}%
    \hfill{\scriptsize\shortbuttons}
  \end{beamercolorbox}}%
  \hbox{\begin{beamercolorbox}[wd=\paperwidth,ht=2.5ex,dp=1ex,leftskip=1em,rightskip=1em]{footline}%
    \insertshortauthor{}\hfill
    \ifshortappendix Appendix\else\insertframenumber/\shortmaincount\fi
  \end{beamercolorbox}}}%
}


\begin{document}

% Original numbered slide 0
\gdef\shortbuttons{}
\begin{frame}[plain,noframenumbering,label=short-0]
  \titlepage
\end{frame}


\section{Introduction}

% Original numbered slide 1
\gdef\shortbuttons{}
\begin{frame}[label=short-1]{Motivation}
  \begin{itemize}
    \item Firms have more access than ever to consumer information
    \item Allows for targeted discounting based on preferences derived from repeat interactions
    \item Targeted discounts allow firms to steer consumers toward particular products.
  \end{itemize}
  \begin{tcolorbox}
    How do firms price based on repeated interactions with consumers and how are profits affected?
  \end{tcolorbox}
\end{frame}


% Original numbered slide 2
\gdef\shortbuttons{}
\begin{frame}[label=short-2]{Research Question}
  \begin{tcolorbox}
    How can firms price discriminate via consumer preferences for novelty and habit and how much is information about these preferences worth?
  \end{tcolorbox}
  How do I answer?
  \begin{itemize}
    \item Solve a theoretical model to understand firm's behavior
    \item Estimate consumer's preferences with Nielsen data, recovering the distribution of novelty/habit preferences
    \item Use the estimated distribution in the pricing model to understand how firm's revenues are impacted by this sort of targeting
  \end{itemize}
\end{frame}


% Original numbered slide 3
\gdef\shortbuttons{}
\begin{frame}[label=short-3]{Literature}
  \begin{columns}[T]
    \begin{column}{0.48\textwidth}
      \textbf{Targeting:}
      \begin{itemize}
        \item Dub\'e (2008)
        \item Dub\'e et. al (2022)
        \item Anand \& Shachar (2009)
        \item Karle \& Reisinger (2024)
        \item Durrmeyer, Fournel, Samano (2026)
      \end{itemize}
    \end{column}
    \begin{column}{0.48\textwidth}
      \textbf{Estimating preferences for novelty:}
      \begin{itemize}
        \item Dub\'e (2004)
        \item Dub\'e et. al (2021)
        \item Chintagunta (1999)
        \item Lattin (1987)
        \item McAllister (1982)
        \item Ackerberg (2003)
      \end{itemize}
    \end{column}
  \end{columns}
\end{frame}


\section{Motivation from Data}

% Original numbered slide 4
\gdef\shortbuttons{\hyperlink{short-5}{\beamergotobutton{Flavor Popularity}}}
\begin{frame}[label=short-4]{HHI}
  \centering
    \includegraphics[width=.95\textwidth,height=.74\textheight,keepaspectratio]{../writeup/1108_writeup/assets/results/summary_statistics/flavor_hhi_histogram.pdf}
\end{frame}


\section{Supply Theory}

% Original numbered slide 6
\gdef\shortbuttons{}
\begin{frame}[label=short-6]{Supply Setting}
  \begin{itemize}
    \item Monopolistic firm producing a single product $j$
    \item Firm is profit maximizing, setting prices according to that objective
    \item A consumer is born into the world with a type $r\in\{H,L\}$, determining how much utility they derive from switching and an initial state $h_0\in\{0,j\}$
    \item Consumers live for two periods, purchasing the inside good or opting for the outside good in each period
    \item The firm knows initial states and types, setting prices to individual consumers by type and history
  \end{itemize}
\end{frame}


% Original numbered slide 7
\gdef\shortbuttons{}
\begin{frame}[label=short-7]{Consumer's Problem}
  \begin{tcolorbox}[
    boxsep=0pt,
    left=2pt, right=2pt,
    top=1pt, bottom=1pt,
    before skip=2pt, after skip=2pt
    ]
    \setlength{\abovedisplayskip}{2pt}
    \setlength{\belowdisplayskip}{2pt}
    \begin{gather*}
      V^{inside}_{jt} = \underbrace{\beta}_{\text{Constant}} + \overbrace{\gamma_r\mathbbm{1}\{h=0\}}^{\text{Utilty from Switching}} - \underbrace{\alpha p_t(h)}_{\text{Price Disutility}}, \ V^{Outside}_{0t} = \underbrace{\gamma_r\mathbbm{1}\{h=j\}}_{\text{Utility from Switching}} \\
      v_t(x) = V_{jt}^{inside} - V_{0t}^{outside} = \beta + \underbrace{\gamma[\mathbbm{1}\{h=0\} - \mathbbm{1}\{h=j\}]}_{q(h)} - \alpha p_t(h)
    \end{gather*}
  \end{tcolorbox}
  \begin{itemize}\setlength{\itemsep}{2pt}
    \item Prices, set by the firm, and switching utility depend on history $h=x_{t-1}$
    \item Utility for the outside good is normalized to zero
  \end{itemize}
\end{frame}


% Original numbered slide 8
\gdef\shortbuttons{}
\begin{frame}[label=short-8]{Firm's Problem}
  \begin{tcolorbox}
    [
    boxsep=0pt,
    left=2pt, right=2pt,
    top=1pt, bottom=1pt,
    before skip=2pt, after skip=2pt
    ]
    \setlength{\abovedisplayskip}{2pt}
    \setlength{\belowdisplayskip}{2pt}
    \begin{gather*}
      \pi_t(h) = \max_{p_t(h)} s_t(h)[p_t(h) - c] \\
      \Pi(h) = \underbrace{s_1(h)(p_1(h) - c)}_{\text{Margin in Period 1}} + \underbrace{\delta\{s_1(h)\pi_2(j) + (1-s_1(h))\pi_2(0)\}}_\text{{Continuation Value}}
    \end{gather*}
  \end{tcolorbox}
  \begin{itemize}
    \item Firm profits are comprised of period one margins, given history of consumption, and the value of that consumer in the different states in period two
    \item The firm will set prices according to consumer type and history to maximize expected profit
  \end{itemize}
\end{frame}


% Original numbered slide 10
\gdef\shortbuttons{\hyperlink{short-9}{\beamergotobutton{Terminal FOC}}\hspace{.35em}\hyperlink{short-11}{\beamergotobutton{Comparative statics}}\hspace{.35em}\hyperlink{short-52}{\beamergotobutton{Elasticities}}}
\begin{frame}[label=short-10]{First Period Optimization}
  \begin{tcolorbox}
    [
    boxsep=0pt,
    left=2pt, right=2pt,
    top=1pt, bottom=1pt,
    before skip=2pt, after skip=2pt
    ]
    \setlength{\abovedisplayskip}{2pt}
    \setlength{\belowdisplayskip}{2pt}
    \begin{align*}
      \frac{d\Pi}{d p_1(h_0)} &= s_1 + \frac{\partial}{\partial p_1(h)}s_1(h_0)[p_1(h_0) - c + \delta(\pi_2(j) - \pi_2(0))] \\
      p_1^*(h_0) &= c + \frac{1}{\alpha(1-s_1^*(h_0))} - \delta(\pi_2(j) - \pi_2(0)) \\
      \Pi^*(h) &= \delta\pi_2(0) + \frac{s_1^*(h_0)}{\alpha(1-s_1^*(h_0))}
    \end{align*}
  \end{tcolorbox}
\end{frame}


% Original numbered slide 14
\gdef\shortbuttons{\hyperlink{short-12}{\beamergotobutton{Calibration}}\hspace{.35em}\hyperlink{short-13}{\beamergotobutton{Table}}\hspace{.35em}\hyperlink{short-15}{\beamergotobutton{Profit}}\hspace{.35em}\hyperlink{short-16}{\beamergotobutton{Price ratio}}\hspace{.35em}\hyperlink{short-17}{\beamergotobutton{Discussion}}}
\begin{frame}[label=short-14]{Optimal Prices}
  \centering
  \includegraphics[width=0.85\textwidth]{assets/theory_delta1/optimal_prices_by_gamma.pdf}
\end{frame}


% Theory plots regenerated at delta=1.
\gdef\shortbuttons{\hyperlink{short-14}{\beamergotobutton{Optimal prices}}\hspace{.35em}\hyperlink{short-15}{\beamergotobutton{Profit detail}}\hspace{.35em}\hyperlink{short-16}{\beamergotobutton{Price-ratio detail}}}
\begin{frame}[label=short-prices-profit]{Optimal Price Ratios and Profit}
  \centering
  \begin{columns}[T,onlytextwidth]
    \begin{column}{.49\textwidth}
      \includegraphics[width=\linewidth,height=.66\textheight,keepaspectratio]{assets/theory_delta1/all_state_price_ratios.pdf}
    \end{column}
    \begin{column}{.49\textwidth}
      \includegraphics[width=\linewidth,height=.66\textheight,keepaspectratio]{assets/theory_delta1/all_state_realized_profit.pdf}
    \end{column}
  \end{columns}
  \vspace{.3em}
  {\footnotesize $\delta=1$: no discounting. Flavor is $j$; outside is everything except that flavor.}
\end{frame}


% Original numbered slide 18
\gdef\shortbuttons{}
\begin{frame}[label=short-18]{Two Pricing Regimes: The Same Consumer Types}
  \small
  \begin{columns}[T,onlytextwidth]
    \begin{column}{.48\textwidth}
      \textbf{Household-specific: type observed}
      \begin{itemize}
        \item Observe type $\gamma_r$ and initial history $h_0$.
        \item Choose $p_{1r}(h_0)$ and then $p_{2r}(h_1)$.
      \end{itemize}
    \end{column}
    \begin{column}{.48\textwidth}
      \textbf{Common: type hidden}
      \begin{itemize}
        \item Observe history and know type weights $\mu_r$.
        \item One offer across types with the same history.
        \item Weights stay fixed; no Bayesian learning.
      \end{itemize}
    \end{column}
  \end{columns}
  \medskip
  \begin{tcolorbox}
    Consumers keep their type in both periods. Evaluate both policies on these
    same persistent types; the gain is household-specific minus common revenue.
  \end{tcolorbox}
  {\footnotesize $\gamma_r\in\{-1.25,0,1.25\}$; equal weights $1/3$ or permutations of $(.50,.25,.25)$.
  In all cases, $h_0$ and subsequent purchases are observed; $\delta=1$, $c=0$.}
\end{frame}


% Original numbered slide 20
\gdef\shortbuttons{\hyperlink{short-19}{\beamergotobutton{Profit levels}}\hspace{.35em}\hyperlink{short-53}{\beamergotobutton{Price paths}}\hspace{.35em}\hyperlink{short-54}{\beamergotobutton{Revenue evaluation}}}
\begin{frame}[label=short-20]{Theory: Revenue Gain from Observing Type}
  \centering
  \includegraphics[width=.98\textwidth,height=.66\textheight,keepaspectratio]{assets/theory_delta1/gamma_uncertainty/information_value.pdf}
  \par\smallskip\footnotesize Expected revenue per consumer, evaluated under the same persistent types.
  Initially outside means outside the focal flavor; it is an initial state, not a population share.
\end{frame}


\section{Data}

% Original numbered slide 21
\gdef\shortbuttons{\hyperlink{short-22}{\beamergotobutton{Activity}}}
\begin{frame}[label=short-21]{Data Sources}
  \textbf{Nielsen HomeScan}
  \begin{itemize}
    \item Provides household level purchase data
    \item Covers years 2006-2015
    \item Products chosen, quantity chosen, timing of trips
  \end{itemize}
  \textbf{Nielsen RetailScan}
  \begin{itemize}
    \item Provides store level product data
    \item For specified product category: prices, markets, product attributes
  \end{itemize}
  \textbf{Sample Description}
  \begin{itemize}
    \item Must be single-person household
    \item Must be seen at least 2 times
    \item Purchases must be made at Nielsen store (prices available)
    \item Flavor attribute must be present
  \end{itemize}
\end{frame}


% Original numbered slide 23
\gdef\shortbuttons{}
\begin{frame}[label=short-23]{Sample Statistics: Demographics}
  \centering
  \begin{adjustbox}{max totalsize={\textwidth}{.83\textheight},center}
    \begin{minipage}{\textwidth}
      \input{../writeup/1108_writeup/assets/results/summary_statistics/full_panel_vs_yogurt.tex}
    \end{minipage}
  \end{adjustbox}
\end{frame}

% Original numbered slide 24
\gdef\shortbuttons{}
\begin{frame}[label=short-24]{Switching Stats}
    \centering
    \begin{adjustbox}{max totalsize={\textwidth}{.75\textheight},center}
    \begin{minipage}{\textwidth}
      \input{../RESULTS_ONLY/summary_statistics/aggregated_flavors/yogurt_dynamics.tex}
    \end{minipage}%
    \end{adjustbox}
\end{frame}


% Original numbered slide 25
\gdef\shortbuttons{}
\begin{frame}[label=short-25]{Persistence Across Aggregated Flavor Groups}
  \centering
  \includegraphics[height=0.69\textheight]{../writeup/1108_writeup/assets/results/summary_statistics/aggregated_flavors/flavor_transitions.pdf}
  \par\small Plain repeats in approximately 80\% of eligible transitions, versus 53\% for berry.
\end{frame}


\section{Demand Estimation}

% Original numbered slide 26
\gdef\shortbuttons{}
\begin{frame}[label=short-26]{Estimate Heterogeneity in Price, Novelty and Habit}
  \begin{itemize}
    \item Weekly choices: other yogurt, berry, plain, or nonpurchase.
    \item Households differ in price sensitivity and responses to recent and
    longer purchase histories.
    \item Flavor intercepts are common across households and differ across flavors.
  \end{itemize}
  \begin{tcolorbox}
    Recover the distribution of switching preferences, then use it in the
    two-period pricing model.
  \end{tcolorbox}
\end{frame}


% Original numbered slide 27
\gdef\shortbuttons{\hyperlink{short-average-history}{\beamergotobutton{Average-history specification}}}
\begin{frame}[label=short-27]{Demand Model}
  \small
  \[
    U_{ijt}=\beta_0+\beta_jX_{jt}+\gamma_iN_{ijt}+\phi_iH_{ijt}
             +\alpha_i p_{jt}+\varepsilon_{ijt},
    \qquad U_{i0t}=\varepsilon_{i0t}.
  \]

  Household $i$'s share of yogurt purchases in flavor $j$ in week $t$ is
  \[
    w_{ijt}=\frac{y_{ijt}}{\sum_k y_{ikt}}.
  \]

  \begin{columns}[T,onlytextwidth]
    \begin{column}{.58\textwidth}
      \textbf{Novelty}
      \[
        N_{ijt}=\begin{cases}
          1-w_{ij,t-1}, & \text{if yogurt bought in $t-1$},\\
          0, & \text{otherwise}.
        \end{cases}
      \]
    \end{column}
    \begin{column}{.35\textwidth}
      \textbf{Habit: fitted specification}
      \[
        H_{ijt}=\sum_{s<t}w_{ijs}.
      \]
    \end{column}
  \end{columns}
\end{frame}

\gdef\shortbuttons{}
\begin{frame}[label=short-identification]{What Variation Separates Tastes from Purchase History?}
  \begin{itemize}
    \item \textbf{Persistent tastes:} repeated choices help estimate stable
    differences across households.
    \item \textbf{Recent history:} changes in last week's basket change the
    novelty predictor within a household.
    \item \textbf{Longer history:} the same last-week basket can follow very
    different earlier purchase patterns.
  \end{itemize}
  \begin{tcolorbox}
    Both histories bought only plain last week and contain 10 earlier purchase weeks.
    Plain's earlier share averages 80\% versus 20\%: novelty is identical,
    average habit is 0.8 versus 0.2.
  \end{tcolorbox}
  \par\smallskip\footnotesize This describes the variation used by the model.
  Separating causal history effects from stable tastes also relies on the
  distributional, initial-history and transitory-shock assumptions.
\end{frame}


% Original numbered slide 29
\gdef\shortbuttons{\hyperlink{short-28}{\beamergotobutton{Probabilities}}\hspace{.35em}\hyperlink{short-55}{\beamergotobutton{Price controls}}\hspace{.35em}\hyperlink{short-56}{\beamergotobutton{Controlled estimates}}\hspace{.35em}\hyperlink{short-61}{\beamergotobutton{First stages}}}
\begin{frame}[label=short-29]{Estimation}
  \begin{tcolorbox}
    [
    boxsep=0pt,
    left=2pt, right=2pt,
    top=1pt, bottom=1pt,
    before skip=2pt, after skip=2pt
    ]
    \setlength{\abovedisplayskip}{2pt}
    \setlength{\belowdisplayskip}{2pt}
    \begin{align*}
      P_i(\theta) &= \int P(y_i\mid\nu,\theta)\ \phi(\nu)\,d\nu \\
      \widehat P_i(\theta) &= \frac{1}{128}\sum_{r=1}^{128}\prod_{t=1}^{T_i}P_(y_i\mid\nu_{ir}, \theta)\qquad\nu_{ir}\sim\mathcal N(0, I) \\
      Q_{128}(\theta) &= -\sum_{i=1}^N\log\widehat P_i(\theta)
    \end{align*}
  \end{tcolorbox}
  \begin{itemize}
    \item $\nu_i$ captures unobserved heterogreneity, $y_{it}$ represents observed choice, $\theta = \{\beta, \gamma, \phi, \alpha\}$ is the vector of heterogeneous demand parameters
    \item 128 draws per household from a normal distribution
    \item 3 starts at different starting values, keep lowest objective
    \item Minimize negative log-likelihood
  \end{itemize}
\end{frame}


% Original numbered slide 32
\gdef\shortbuttons{\hyperlink{short-30}{\beamergotobutton{Estimates}}\hspace{.35em}\hyperlink{short-33}{\beamergotobutton{Discussion}}}
\begin{frame}[label=short-32]{Distributions}
  \includegraphics[width=.92\textwidth,height=.50\textheight,keepaspectratio]{assets/demand_1111_three/history_distributions.pdf}
\end{frame}


% Original numbered slide 31
\gdef\shortbuttons{\hyperlink{short-62}{\beamergotobutton{Controlled WTP}}}
\begin{frame}[label=short-31]{Implied Willingness to Pay}
  \centering
  \input{assets/demand_1111_three/implied_wtp.tex}
  \begin{itemize}
    \item Would need to be subsidized to switch, though lots of variation
    \item Willing to pay to continue consuming their same product
  \end{itemize}
\end{frame}


\section{Pricing Exercise}

% Original numbered slide 34
\gdef\shortbuttons{}
\begin{frame}[label=short-34]{Quantifying the Theory with Estimated Preferences}
  \[
    v_t=\beta_f+\alpha p_t+\gamma\big[\mathbbm{1}\{h_t=0\}-\mathbbm{1}\{h_t=j\}\big].
  \]
  \begin{itemize}
    \item Same binary choices and pricing problem as the theory: plain ($j$) versus not plain ($0$).
    \item Use the estimated novelty distribution, mean price sensitivity and plain intercept.
    \item Hold habit at its reference level; no additional history dynamics.
  \end{itemize}
  \begin{tcolorbox}
    Compare informed and common prices, and the revenue each generates.
    This is a quantified theory illustration using saved estimates.
  \end{tcolorbox}
\end{frame}

\gdef\shortbuttons{}
\begin{frame}[label=short-40]{Only Knowledge of Novelty Changes}
  \begin{itemize}
    \item Informed: the firm observes the consumer's $\gamma$.
    \item Common: the firm knows the distribution of $\gamma$ and offers one price across types.
    \item Both observe purchase history and adapt the second price to plain versus not plain.
    \item Common-price beliefs remain fixed, as in the theory; consumer types persist.
  \end{itemize}
  \begin{tcolorbox}
    Evaluate both policies on the same consumers over two periods.
  \end{tcolorbox}
\end{frame}

% Revised empirical information comparison; populated by aggregate-only refresh.
\IfFileExists{assets/quantified_theory/plain/optimal_prices_by_gamma.pdf}{%
\gdef\shortbuttons{}
\begin{frame}[label=short-revised-prices]{Quantified Theory: Optimal Prices across Novelty Types}
  \centering
  \includegraphics[width=.98\textwidth,height=.68\textheight,keepaspectratio]{assets/quantified_theory/plain/optimal_prices_by_gamma.pdf}
  \par\smallskip\footnotesize Same model and layout as supply, using estimated preference parameters.
  There are exactly two terminal prices, after plain and after not plain.
\end{frame}

\gdef\shortbuttons{}


\gdef\shortbuttons{}
\begin{frame}[label=short-revised-paths]{Quantified Theory: Informed and Common Price Paths}
  \centering
  \includegraphics[width=.98\textwidth,height=.68\textheight,keepaspectratio]{assets/quantified_theory/plain/information_price_paths_comparison.pdf}
  \par\smallskip\footnotesize Solid paths average informed prices within novelty groups; dashed paths show common prices.
  Both regimes choose the terminal price after observing the first-period binary choice.
\end{frame}

\gdef\shortbuttons{}


\gdef\shortbuttons{}
\begin{frame}[label=short-revised-revenue]{Revenue under the Same Persistent Novelty Types}
  \centering
  \includegraphics[width=.98\textwidth,height=.68\textheight,keepaspectratio]{assets/quantified_theory/plain/information_revenues.pdf}
  \par\smallskip\footnotesize Expected two-period revenue per model consumer; zero marginal cost.
  Common prices are history-adaptive, with novelty hidden and fixed beliefs.
\end{frame}

\gdef\shortbuttons{}
\begin{frame}[label=short-revised-gain]{How Much Revenue Does Knowing Novelty Add?}
  \centering
  \includegraphics[width=.98\textwidth,height=.68\textheight,keepaspectratio]{assets/quantified_theory/plain/information_value.pdf}
  \par\smallskip\footnotesize Informed minus common revenue, evaluated on the same persistent types.
  Prior sensitivities use one novelty group with probability one-half.
\end{frame}

\gdef\shortbuttons{}
\begin{frame}[label=short-revised-table]{Revenue Comparison by Initial History and Prior}
  \centering\small
  \begin{adjustbox}{max width=\linewidth}
    \input{assets/quantified_theory/plain/revenue_table.tex}
  \end{adjustbox}
  \par\smallskip\footnotesize Dollars per model consumer over two periods.
  Both regimes adapt second prices to history; only novelty information differs.
\end{frame}
}{%
\gdef\shortbuttons{}
\begin{frame}[label=short-revised-pending]{Binary Plain-versus-Not Pricing: Rerun Pending}
  \begin{itemize}
    \item Saved estimates are inputs to the binary model, including known habit.
    \item Recompute household-specific and common prices over two binary purchase periods.
    \item Only novelty is hidden; price sensitivity and habit are known.
    \item Report price changes and revenue gains in the same figure layouts as supply.
  \end{itemize}
  \begin{tcolorbox}
    Revised results require a pricing-only server run.
    Earlier full-basket results do not apply to this binary specification.
  \end{tcolorbox}
\end{frame}
}





% Original numbered slide 46
\gdef\shortbuttons{}
\begin{frame}[label=short-46]{Calibrated Plain Illustration: Revenue Levels}
  \centering
  \includegraphics[width=.98\textwidth,height=.67\textheight,keepaspectratio]{assets/aligned_pricing/information_friction_profits.pdf}
  \par\smallskip\footnotesize Dollars per model consumer over two periods.
  Both policies are evaluated under the same persistent types; $\delta=1$.
\end{frame}





% Original numbered slide 47


% Original numbered slide 57






\section{Conclusion}

% Original numbered slide 50
\gdef\shortbuttons{}
\begin{frame}[label=short-50]{What Does Knowing Novelty Change?}
  \begin{itemize}
    \item Estimated preferences quantify the theoretical information comparison.
    \item Compare acquisition and repeat-purchase prices separately for initially inside and outside consumers.
    \item Measure the revenue gain from knowing novelty on the same persistent consumer types.
  \end{itemize}
  \begin{tcolorbox}
    \IfFileExists{assets/quantified_theory/plain/gain_headline.tex}{%
      \input{assets/quantified_theory/plain/gain_headline.tex}
    }{Compare informed and common revenue under the estimated novelty distribution.}
  \end{tcolorbox}
  \par\smallskip\footnotesize Fitted-distribution baseline; expected revenue per model consumer over two periods.
  Both regimes adapt prices to purchase history; only knowledge of novelty differs.
\end{frame}


% Original numbered slide 51



\appendix

\shortappendixtrue

\section{Appendix}

% Original numbered slide 22
\gdef\shortbuttons{\hyperlink{short-21}{\beamergotobutton{Back to Data Sources}}}
\begin{frame}[label=short-22,noframenumbering]{Sample Statistics: Purchase Activity}
  \centering
  \small
  \renewcommand{\arraystretch}{1.15}
  \begin{tabular}{lrr}
    \toprule
     & Full sample & Yogurt sample \\
    \midrule
    Households & 37,827 & 15,106 \\
    Observed panel years & 3.96 & 3.16 \\
    Trips per household & 493.13 & 20.26 \\
    Trips per observed week (approx.) & 2.27 & 0.11 \\
    Distinct UPCs per household & 1,178.31 & 13.89 \\
    Purchase lines per household & 3,397.64 & 45.15 \\
    Packages per household & 4,669.07 & 93.77 \\
    \bottomrule
  \end{tabular}

  \medskip
  {\footnotesize All statistics except household counts are means.}
\end{frame}





% Original numbered slide 5
\gdef\shortbuttons{\hyperlink{short-4}{\beamergotobutton{Back to HHI}}}
\begin{frame}[label=short-5,noframenumbering]{Flavor Popularity}
  \centering
    \includegraphics[width=.95\textwidth,height=.74\textheight,keepaspectratio]{../writeup/1108_writeup/assets/results/summary_statistics/flavor_shares.pdf}
\end{frame}


% Original numbered slide 9
\gdef\shortbuttons{\hyperlink{short-10}{\beamergotobutton{Back to First Period Optimization}}}
\begin{frame}[label=short-9,noframenumbering]{Optimal Terminal Price and Profit}
  \begin{tcolorbox}
    [
    boxsep=0pt,
    left=2pt, right=2pt,
    top=1pt, bottom=1pt,
    before skip=2pt, after skip=2pt
    ]
    \setlength{\abovedisplayskip}{2pt}
    \setlength{\belowdisplayskip}{2pt}
    \begin{align*}
      \frac{\partial\pi_2(h_1)}{\partial p_2(h_1)} &= s_2(h_1) + (p_2 - c)\frac{\partial s_2(h_1)}{\partial p_2(h_1)} \\
      p_2^*(h_1)&= c + \frac{1}{\alpha(1-s_2^*(h_1))} \\
      \pi^*_2(p_2^*(h_1); h_1) &= \frac{s_2^*(h_1)}{\alpha(1-s_2^*(h_1))}
    \end{align*}
  \end{tcolorbox}
  \begin{itemize}
    \item Optimal period two price is cost plus a demand markup term
    \item Period two profit at optimal prices is the odds of purchase, scaled by the inverse price sensitivity
  \end{itemize}
\end{frame}


% Original numbered slide 11
\gdef\shortbuttons{\hyperlink{short-10}{\beamergotobutton{Back to First Period Optimization}}}
\begin{frame}[label=short-11,noframenumbering]{Optimal Profit Given Switching Type}
  \begin{tcolorbox}
    [
    boxsep=0pt,
    left=2pt, right=2pt,
    top=1pt, bottom=1pt,
    before skip=2pt, after skip=2pt
    ]
    \setlength{\abovedisplayskip}{2pt}
    \setlength{\belowdisplayskip}{2pt}
    \begin{align*}
      \frac{d\Pi}{d\gamma} &= \frac{q(h)s_1^*(h) - \delta s_1^*(h)s_2^*(j) + \delta(1-s_1^*(h))s_2^*(0)}{\alpha}
    \end{align*}
  \end{tcolorbox}
  \begin{itemize}
    \item There are two interpretations:
    \begin{enumerate}
      \item If $h_0=j$, $q(h)<0$, leading us to ambiguous signage as we raise profits from initially outside consumers but see opposing forces for the initially inside: leave in period one, return in period two
      \item If $h_0=0$, $q(h)>0\Rightarrow \frac{d\Pi}{d\gamma} > 0$, meaning that profits increase as $\gamma$ increases
    \end{enumerate}
    \item Points to firms potentially reaping great benefit from knowing consumer types and histories
  \end{itemize}
\end{frame}


% Original numbered slide 52
\gdef\shortbuttons{\hyperlink{short-10}{\beamergotobutton{Back to First Period Optimization}}}
\begin{frame}[label=short-52,noframenumbering]{Shares and Elasticities}
  \begin{tcolorbox}
    [
    boxsep=0pt,
    left=2pt, right=2pt,
    top=1pt, bottom=1pt,
    before skip=2pt, after skip=2pt
    ]
    \setlength{\abovedisplayskip}{2pt}
    \setlength{\belowdisplayskip}{2pt}
    \begin{gather*}
      s_t(h) = \frac{e^{v_t(h)}}{1+e^{v_t(h)}}, \ s_{0t}(h) = 1-s_t(h) \\
      \frac{\partial s_t(h)}{\partial p_t(h)} = -\alpha s_t(h)(1-s_t(h)) \\
      \frac{\partial s_t(h)}{\partial\gamma} = q(h)s_t(h)(1-s_t(h))
    \end{gather*}
  \end{tcolorbox}
  \end{frame}


% Original numbered slide 12
\gdef\shortbuttons{\hyperlink{short-14}{\beamergotobutton{Back to Optimal Prices}}}
\begin{frame}[label=short-12,noframenumbering]{Pricing Simulations Specification}
  \centering
  \small
  \begin{tabular}{lll}
    \toprule
    Parameter & Description & Value \\
    \midrule
    $\beta$ & Inside-good intercept & $2$ \\
    $\alpha$ & Price sensitivity & $1$ \\
    $\delta$ & Discount factor (no discounting) & $1$ \\
    $c$ & Marginal cost & $0$ \\
    $\gamma$ & Switching preference & $[-3,3]$ \\
    \midrule
    \multicolumn{2}{l}{Initial Consumption State}
      & $h_0\in\{0,j\}$ \\
    \multicolumn{2}{l}{Feasible Prices}
      & $p_t(h)\in\mathbb{R}$ \\
    \bottomrule
  \end{tabular}
  \begin{itemize}
    \item Allow for negative pricing
    \item Firm has knowledge of both type and history, allowing for targeted pricing
    \item Four consumption histories possible $(h_0, h_1) \in \{(0,0), (0,j), (j,j), (j,0)\}$
  \end{itemize}
\end{frame}


% Original numbered slide 13
\gdef\shortbuttons{\hyperlink{short-14}{\beamergotobutton{Back to Optimal Prices}}}
\begin{frame}[label=short-13,noframenumbering]{Results}
  \begin{center}\small\begin{tabular}{crrrrrr}\toprule
  $x_0$ & $\gamma$ & $p_1^*$ & $p_2^*(j)$ & $p_2^*(0)$ & $\Pi^*$ & $\Pi^{*\prime}$\\\midrule
  $j$ & -3 & 3.358 & 3.926 & 1.120 & 5.284 & -1.445\\
  $j$ & -1 & 2.211 & 2.557 & 1.567 & 2.768 & -0.993\\
  $j$ & 0 & 2.000 & 2.000 & 2.000 & 2.000 & -0.500\\
  $j$ & 1 & 2.271 & 1.567 & 2.557 & 1.838 & 0.177\\
  $j$ & 3 & 3.814 & 1.120 & 3.926 & 2.934 & 0.730\\
  $0$ & -3 & -0.901 & 3.926 & 1.120 & 1.026 & 0.177\\
  $0$ & -1 & 1.005 & 2.557 & 1.567 & 1.562 & 0.376\\
  $0$ & 0 & 2.000 & 2.000 & 2.000 & 2.000 & 0.500\\
  $0$ & 1 & 2.995 & 1.567 & 2.557 & 2.562 & 0.624\\
  $0$ & 3 & 4.905 & 1.120 & 3.926 & 4.025 & 0.823\\
  \bottomrule\end{tabular}\end{center}
\end{frame}


% Original numbered slide 15
\gdef\shortbuttons{\hyperlink{short-14}{\beamergotobutton{Back to Optimal Prices}}}
\begin{frame}[label=short-15,noframenumbering]{Optimal Profit}
  \centering
  \includegraphics[width=.75\textwidth,height=.68\textheight,keepaspectratio]{assets/theory_delta1/all_state_realized_profit.pdf}
\end{frame}


% Original numbered slide 16
\gdef\shortbuttons{\hyperlink{short-14}{\beamergotobutton{Back to Optimal Prices}}}
\begin{frame}[label=short-16,noframenumbering]{Optimal Price Ratio}
  \centering
  \includegraphics[width=.75\textwidth,height=.68\textheight,keepaspectratio]{assets/theory_delta1/all_state_price_ratios.pdf}
\end{frame}


% Original numbered slide 17
\gdef\shortbuttons{\hyperlink{short-14}{\beamergotobutton{Back to Optimal Prices}}}
\begin{frame}[label=short-17,noframenumbering]{Discussion}
  \begin{itemize}
    \item Two mechanisms: Induce and harvest, Retain and gain
    \item Most profit from those who are habitual inside purchasers
    \item Least profit from those who purchase outside in period 1 across all levels of gamma
    \item Interesting when $(j,j)$ and $(0,0)$ cross
  \end{itemize}
\end{frame}


% Original numbered slide 19
\gdef\shortbuttons{\hyperlink{short-20}{\beamergotobutton{Back to Household-Specific vs Common Prices}}}
\begin{frame}[label=short-19,noframenumbering]{Theory: Revenue Levels under the Same Persistent Types}
  \centering
  \includegraphics[width=.98\textwidth,height=.72\textheight,keepaspectratio]{assets/theory_delta1/gamma_uncertainty/profits.pdf}
  \par\smallskip\footnotesize Full-information optimal revenue versus realized revenue of the fixed-belief common policy.
\end{frame}


% Original numbered slide 53
\gdef\shortbuttons{\hyperlink{short-20}{\beamergotobutton{Back to Household-Specific vs Common Prices}}}
\begin{frame}[label=short-53,noframenumbering]{Theory: Type-Specific and Common Offers}
  \centering
  \includegraphics[width=.98\textwidth,height=.70\textheight,keepaspectratio]{assets/theory_delta1/gamma_uncertainty/information_price_paths_comparison.pdf}
  \par\smallskip\footnotesize Equal type weights; same format as the calibrated plain price paths.
\end{frame}


% Original numbered slide 54
\gdef\shortbuttons{\hyperlink{short-20}{\beamergotobutton{Back to Household-Specific vs Common Prices}}}
\begin{frame}[label=short-54,noframenumbering]{Fixed Beliefs for Offers; Persistent Types for Evaluation}
  \begin{itemize}
    \item The hidden-type firm uses $\mu_r$ in every state to choose offers;
    it does not update beliefs after a purchase.
    \item Real consumers retain their type. First-period choices sort types into
    histories even when the firm ignores that information.
  \end{itemize}
  \[
    \Pr(j\text{ in both periods})=\sum_r\mu_r s_{1r}s_{2r}(j),
    \quad\text{not}\quad
    \left(\sum_r\mu_r s_{1r}\right)\left(\sum_r\mu_r s_{2r}(j)\right).
  \]
  \begin{tcolorbox}
    The firm's fixed-belief optimization objective and realized persistent-type
    revenue need not coincide. All displayed policy comparisons use the latter.
  \end{tcolorbox}
\end{frame}


% Original numbered slide 28
\gdef\shortbuttons{\hyperlink{short-29}{\beamergotobutton{Back to Estimation}}}
\begin{frame}[label=short-28,noframenumbering]{Probability and Shares}
  \begin{align*}
    P_{ijt} &= \frac{\text{exp}\{V_{ijt}\}}{1 + \sum_{k\in J}\text{exp}\{V_{ikt}\}} \\
    &\text{s.t.} \ V_{ijt} = U_{ijt} - \varepsilon_{ijt} \\
    s_{jt} &= \frac{1}{N}\sum_{n=1}^{N}P_{ijt}
  \end{align*}
  \begin{itemize}
    \item Use deterministic utility to form purchase probabilities
    \item Take average probability for product $j$ across individuals $i\in N$ to recover shares
  \end{itemize}
\end{frame}


% Original numbered slide 55
\gdef\shortbuttons{\hyperlink{short-29}{\beamergotobutton{Back to Estimation}}}
\begin{frame}[label=short-55,noframenumbering]{Demand: Price-Control Comparison}
\input{../writeup/1108_writeup/assets/results/demand/comparison.tex}
\end{frame}


% Original numbered slide 56
\gdef\shortbuttons{\hyperlink{short-29}{\beamergotobutton{Back to Estimation}}}
\begin{frame}[label=short-56,noframenumbering]{Controlled Demand Estimates}
\input{../writeup/1108_writeup/assets/results/demand_controlled/estimates.tex}
\end{frame}


% Original numbered slide 61
\gdef\shortbuttons{\hyperlink{short-29}{\beamergotobutton{Back to Estimation}}}
\begin{frame}[label=short-61,noframenumbering]{Controlled Fit: Price First Stages}
\input{../writeup/1108_writeup/assets/results/demand_controlled/first_stage.tex}
\end{frame}


% Original numbered slide 30
\gdef\shortbuttons{\hyperlink{short-32}{\beamergotobutton{Back to Distributions}}}
\begin{frame}[label=short-30,noframenumbering]{Estimates}
  \centering
  \input{assets/demand_1111_three/estimates.tex}
\end{frame}


% Original numbered slide 33
\gdef\shortbuttons{\hyperlink{short-32}{\beamergotobutton{Back to Distributions}}}
\begin{frame}[label=short-33,noframenumbering]{Discussion}
  \begin{itemize}
    \item Mean novelty WTP is approximately $-\$0.84$
    \item Mean habit preference is positive
    \item Berry is less preferred than other; plain is similar to other
  \end{itemize}
  \begin{tcolorbox}
    How can the firm use the distribution of switching types to price discriminate?
  \end{tcolorbox}
\end{frame}


% Original numbered slide 62
\gdef\shortbuttons{\hyperlink{short-31}{\beamergotobutton{Back to Implied Willingness to Pay}}}
\begin{frame}[label=short-62,noframenumbering]{Controlled Fit: Willingness to Pay}
\input{../writeup/1108_writeup/assets/results/demand_controlled/implied_wtp.tex}
\end{frame}


% Original numbered slide 37
\gdef\shortbuttons{\hyperlink{short-44}{\beamergotobutton{Back to Revenue Gain}}}
\begin{frame}[label=short-37,noframenumbering]{Results: Plain Inside}
  \centering
  \includegraphics[width=\textwidth]{../../../Results/pricing_no_disc/pricing/plain/household_coupon_results.pdf}
\end{frame}


% Original numbered slide 39
\gdef\shortbuttons{\hyperlink{short-44}{\beamergotobutton{Back to Full-Demand Prices}}}
\begin{frame}[label=short-39,noframenumbering]{Plain: Optimal Offers by Consumer Type}
  \centering
  \includegraphics[width=0.85\textwidth]{../../../Results/pricing_no_disc/pricing/plain/discounts_by_consumer_type.pdf}
\end{frame}




\gdef\shortbuttons{\hyperlink{short-plain-complete-paths}{\beamergotobutton{Back to Calibrated Prices}}}
\begin{frame}[label=short-plain-gamma-detail,noframenumbering]{Calibrated Plain: Detail around Estimated Group Means}
  \centering
  \includegraphics[width=.98\textwidth,height=.69\textheight,keepaspectratio]{assets/aligned_pricing/calibrated_prices_by_gamma_detail.pdf}
  \par\smallskip\footnotesize Same curves, zoomed around low, middle and high group means.
  All three means are negative. Terminal curves condition on the first-period choice.
\end{frame}



\gdef\shortbuttons{\hyperlink{short-plain-info-table}{\beamergotobutton{Back to Revenue Comparison}}\hspace{.35em}\hyperlink{short-37}{\beamergotobutton{Plain types}}\hspace{.35em}\hyperlink{short-39}{\beamergotobutton{Price offers}}}
\begin{frame}[label=short-44,noframenumbering]{Separate Full-Demand Exercise: Earlier Shared-Schedule Price Benchmark}
  \centering
  \includegraphics[width=.85\textwidth,height=.57\textheight,keepaspectratio]{assets/aligned_pricing/pricing_policy_benchmarks.pdf}
  \par\smallskip\small Plain: \input{assets/aligned_pricing/coupon_headline.tex}.
  \par\smallskip\footnotesize Same dynamic-pricing mechanism, with multinomial demand, novelty and habit.
  Common prices use shared percentage reductions from posted prices;
  household-specific prices use estimated types and basket histories.
  The middle group remains at posted prices. Totals are not per-consumer binary-model revenue.
\end{frame}

\gdef\shortbuttons{\hyperlink{short-34}{\beamergotobutton{Back}}}
\begin{frame}[label=short-35,noframenumbering]{Earlier Binary Illustration: Calibration Details}
  \begin{itemize}
    \item Low, middle and high groups are defined by estimated novelty preferences.
    Use each group's mean; the middle type is \textbf{not} set to zero.
    \[
      (\gamma_L,\gamma_M,\gamma_H)=(-1.112,\,-0.763,\,-0.223).
    \]
    \item These are relative low/high types: all three group means are negative.
    \item Map novelty to the binary model's symmetric history effect,
    $\gamma_r[\mathbbm{1}\{h=0\}-\mathbbm{1}\{h=j\}]$.
    \item Calibrate a common intercept and price sensitivity to observed offered prices.
    Keep $\delta=1$ and $c=0$.
  \end{itemize}
  \par\smallskip\footnotesize Illustration, not the full split-history demand model:
  habit is not separately carried into this binary calibration.
  Price targets use 32 complete histories; treat the calibration as preliminary.
\end{frame}

\gdef\shortbuttons{\hyperlink{short-40}{\beamergotobutton{Back}}}
\begin{frame}[label=short-42,noframenumbering]{Revenue Evaluation with Full Basket Histories}
  \small
  \[
    R(h_0)=\sum_r\mu_r\sum_{b_1\in\mathcal B}
      P_r(b_1\mid p_{1r},h_0)
      \left[p_{1r}b_{1j}
       +\delta\,p_{2r}(b_1)K s_{rj}\big(p_{2r}(b_1),h_1(b_1)\big)\right].
  \]
  \begin{itemize}
    \item $b_1$ is the entire first-period binary choice; $b_{1j}$ counts focal purchases.
    \item $h_1(b_1)$ updates both novelty and habit using the saved demand specification.
    \item Household-specific prices vary with $r$. Common prices do not, conditional
    on the other known characteristics and the same basket history.
    \item The same type $r$ determines purchases in both periods.
  \end{itemize}
  \begin{tcolorbox}
    Limited-information offers use fixed-belief optimization. Reported revenue
    evaluates those offers under persistent types, as in the theory comparison.
  \end{tcolorbox}
\end{frame}

\gdef\shortbuttons{\hyperlink{short-50}{\beamergotobutton{Back}}}
\begin{frame}[label=short-51,noframenumbering]{Next Steps}
  \begin{itemize}
    \item Add firm learning to supply side
    \item Extend market structure on supply
    \item Add more product attributes to empirical estimation
    \item Estimate unobserved quality in BLP fashion
  \end{itemize}
\end{frame}

\gdef\shortbuttons{\hyperlink{short-27}{\beamergotobutton{Back to Demand Model}}}
\begin{frame}[label=short-average-history,noframenumbering]{Average History: Fitted Specification}
  \[
    H^{\mathrm{avg}}_{ijt}=\frac{1}{n_{it}}
      \sum_{s\in\mathcal{O}_i,\,s<t}w_{ijs},
    \qquad n_{it}=\#\{s\in\mathcal{O}_i:s<t\}.
  \]
  \begin{itemize}
    \item $\mathcal{O}_i$: observed yogurt-purchase weeks; $n_{it}$ counts prior occasions.
    \item Empty history gives zero. Unobserved calendar weeks do not enter the denominator.
    \item Current purchases affect the next state only.
  \end{itemize}
  \begin{tcolorbox}
    Displayed demand estimates use the November 11 three-flavor average-history fit.
    The pricing illustration holds habit at its reference level.
  \end{tcolorbox}
\end{frame}

\gdef\shortbuttons{\hyperlink{short-34}{\beamergotobutton{Back to Revised Exercise}}}
\begin{frame}[label=short-plain-complete-paths,noframenumbering]{Earlier Binary Illustration: Optimal Prices across Switching Types}
  \centering
  \includegraphics[width=.98\textwidth,height=.68\textheight,keepaspectratio]{assets/aligned_pricing/calibrated_prices_by_gamma.pdf}
  \par\smallskip\footnotesize Same known-type model as the theory price plots;
  plain calibration: $\beta=10.502$, price sensitivity $\alpha=3.187$, $\delta=1$, $c=0$.
  Points mark estimated group means; shading spans those means.
  The wider $\gamma$ range is a comparative-static exercise, not estimated support.
\end{frame}

\gdef\shortbuttons{\hyperlink{short-34}{\beamergotobutton{Back to Revised Exercise}}}
\begin{frame}[label=short-plain-price-paths,noframenumbering]{Earlier Binary Illustration: Type-Specific and Common Price Paths}
  \centering
  \includegraphics[width=.98\textwidth,height=.54\textheight,keepaspectratio]{assets/aligned_pricing/information_price_paths_comparison.pdf}
  \par\smallskip\footnotesize Equal type weights. Paths show $p_1\rightarrow\mathbb{E}[p_2]$.
  Solid paths distinguish low, middle and high types; the dashed path is the common offer.
  Panels distinguish initially inside and outside consumers. Terminal offers average over first-period choices.
\end{frame}

\gdef\shortbuttons{\hyperlink{short-34}{\beamergotobutton{Back to Revised Exercise}}}
\begin{frame}[label=short-47,noframenumbering]{Earlier Binary Illustration: Revenue Gain from Observing Type}
  \centering
  \includegraphics[width=.98\textwidth,height=.55\textheight,keepaspectratio]{assets/aligned_pricing/information_friction_information_value.pdf}
  \par\smallskip\footnotesize Gain per consumer over two periods, using the same persistent types under both policies.
\end{frame}

\gdef\shortbuttons{\hyperlink{short-34}{\beamergotobutton{Back to Revised Exercise}}}
\begin{frame}[label=short-57,noframenumbering]{Earlier Binary Illustration: Information Gains by Initial State}
  \begin{center}
    \begin{tabular}{lrr}\toprule
      Equal type weights & Initially inside & Initially outside\\\midrule
      Theory: personalization gain & \$0.0555 & \$0.2642\\
      Calibrated plain: personalization gain & \$0.0249 & \$0.0358\\\bottomrule
    \end{tabular}
  \end{center}
  \begin{itemize}
    \item Attracting different types can require very different offers.
    A common offer sacrifices the ability to tailor those incentives.
    \item Information value depends on the separation between optimal type-specific offers,
    not simply the revenue level or initial purchase rate.
    \item In the calibrated example, initially outside consumers generate less total
    revenue but a larger \textbf{incremental gain} from personalization.
  \end{itemize}
\end{frame}

\gdef\shortbuttons{\hyperlink{short-34}{\beamergotobutton{Back to Revised Exercise}}}
\begin{frame}[label=short-plain-info-table,noframenumbering]{Earlier Binary Illustration: Revenue by State and Prior}
  \centering\small
  \begin{adjustbox}{max width=\linewidth}
    \input{assets/aligned_pricing/information_results.tex}
  \end{adjustbox}
  \par\medskip\footnotesize Dollars per model consumer, evaluated under persistent types.
  $j$: initially plain; $0$: initially outside plain, including other yogurt.
  Price targets use 32 complete histories; this is an illustrative calibration.
\end{frame}

\end{document}
